\section{Tetraquarks and pentaquarks: hadrons as trees of dualities}\label{sec:exotics}

Hypothesis~\ref{h:tree} (Section~\ref{sec:matter}) reads Axioms~1 and~3 for the strong force: every hadron is a binary tree of two
colour-conjugate halves, each internal node a bound duality. It was fixed before any analysis, as statement SM6 of the pre-registration
next to the Standard Model statements of Section~\ref{sec:sm}. Hadron spectroscopy is measured to MeV precision, so H8 can be tested where its content lies: in the masses of the compact
trees, in the closeness of molecular states to two-hadron thresholds, and in the partner states a tree implies. We test all three
with one constituent-quark model with a colour-spin interaction, fitted to ordinary hadrons only. The generic statement holds as a
reading, because QCD's colour algebra already allows every hadron to be written this way. The model's compact trees fail more often
than they hold: 4 of 9 scored masses hold, 5 are refuted, and the doubly heavy tetraquarks $T_{bb}$ and $T_{bc}$ remain live
predictions. Exotic hadrons lie closer to two-hadron thresholds than chance placement would put them, but that is the known signature
of hadronic molecules, it rests on few states, and it does not single out GUID.

\subsection{What H8 says beyond QCD}\label{sec:exotics-qcd}

A quark is a colour triplet $\mathbf 3$, an antiquark a $\mathbf{\bar 3}$, and a hadron a colour singlet. The products
$\mathbf3\otimes\mathbf3=\mathbf{\bar3}\oplus\mathbf6$, $\mathbf{\bar3}\otimes\mathbf{\bar3}=\mathbf3\oplus\mathbf{\bar6}$ and
$\mathbf3\otimes\mathbf{\bar3}=\mathbf1\oplus\mathbf8$ each contain the conjugate triplet or the singlet, so any set of quarks and
antiquarks that can form a singlet can be coupled pairwise, node by node, until two conjugate halves meet at the root \tagP. That is
the generic statement of H8, and it holds in QCD for every hadron, observed or not. The two branches of H8 also exhaust the colour
space of a tetraquark: the singlets of $qq\bar q\bar q$ form a two-dimensional space, and the compact tree
$[q_1q_2]_{\bar3}[\bar q_3\bar q_4]_3$ and the molecular tree $(q_1\bar q_3)_1(q_2\bar q_4)_1$ are two vectors of it with overlap
$|\langle\text{compact}|\text{molecular}\rangle|^2=1/3$ \chk{exotics:E}, so every tetraquark colour state is a superposition of the two
H8 trees. With its either-or, H8 can therefore not fail on any single state. \emph{Verdict: holds as a reading.}

The content lies in the clause that every node is a \emph{bound} duality, made quantitative by a model, and in what the trees imply:
(a) masses of the compact trees, (b) closeness to threshold where both halves are colour singlets, (c) isospin and spin partners. We
use no numerology. In particular we do not connect the five quarks of a pentaquark with the (at least) five parts of the minting (G0,
Section~\ref{sec:dmp}): the number of constituents of a hadron is fixed by colour algebra and baryon number, the minting's parts are
holders of root keys, and the axioms supply no map between the two.

\subsection{The model}\label{sec:exotics-model}

\paragraph{The colour-spin Hamiltonian.} We use the constituent-quark model with a chromomagnetic hyperfine interaction of De
R\'ujula, Georgi and Glashow~\cite{derujula1975}, in the form of Karliner and Rosner~\cite{karliner2014,karliner2017} (KR) that predicted
the doubly charmed baryon and the doubly heavy tetraquarks. The mass of a hadron with constituents $i$ is
\begin{equation}\label{eq:exo-H}
M=\sum_i m_i+\sum_{i<j}B_{ij}-\sum_{i<j}C_{ij}\,\langle\boldsymbol\lambda_i\!\cdot\!\boldsymbol\lambda_j\,
\boldsymbol\sigma_i\!\cdot\!\boldsymbol\sigma_j\rangle ,
\end{equation}
with $\boldsymbol\lambda_i\!\cdot\!\boldsymbol\lambda_j=\sum_{a=1}^8\lambda^a_i\lambda^a_j$ the colour product of the Gell-Mann
matrices ($-\lambda^{a*}$ for an antiquark), $\boldsymbol\sigma$ the Pauli matrices, $m_i$ effective masses, $B_{ij}$ binding terms
and $C_{ij}$ hyperfine strengths. The expectation value is taken in the H8 tree: every internal node below the root couples its two
branches to a colour (anti)triplet, the root to a singlet, and every node carries a spin. The colour factors that follow are
$\langle\boldsymbol\lambda_i\!\cdot\!\boldsymbol\lambda_j\rangle=-16/3$ for a quark and antiquark in a singlet, $-8/3$ for a pair in a
diquark (and every pair of a baryon), $-4/3$ for a quark and an antiquark in the two halves of $[qq][\bar q\bar q]$, $-2/3$ between the
two diquarks of $([qq][qq])\bar q$, and $-4/3$ between its antiquark and each quark \chk{exotics:E}. Within a fixed colour tree, node
spins that give the same total $J$ are mixed and the Hamiltonian is diagonalised; for neutral states with hidden flavour the eigenvectors
are sorted by $C$ parity. Two identical $s$, $c$ or $b$ quarks in a diquark have spin 1 (Pauli), and a light $[ud]$ pair has isospin
equal to its spin.

\paragraph{Couplings, as in KR.} Mesons and baryons have separate effective masses $m^m_f$ and $m^b_f$ ($f=q,s,c,b$, with $q=u,d$),
and KR's convention gives every constituent of a compact multiquark the baryon mass. The coefficient of $\boldsymbol\sigma_i\!\cdot\!
\boldsymbol\sigma_j$ is $v^m_{ij}=2b\,k_{ij}/(m^m_im^m_j)$ in a meson and $v^b_{ij}=a\,k_{ij}/(m^b_im^b_j)$ in a baryon, so
\begin{equation}\label{eq:exo-C}
C_{ij}=\tfrac{3}{16}\,v^m_{ij}\ \ (\text{quark--antiquark pair}),\qquad C_{ij}=\tfrac38\,v^b_{ij}\ \ (\text{quark--quark pair}),
\end{equation}
which returns $v^m$ and $v^b$ for the colour factors $-16/3$ and $-8/3$. The factor $k_{ij}$ is 1 except for a $(c,s)$ or $(b,s)$ pair,
where it is KR's enhancement $k_s=a_{cs}/a$; KR use it in baryons, and we share it with mesons, where KR found the same enhancement of
the $D_s^*-D_s$ splitting. Heavy--heavy pairs take $v^m_{c\bar c}=[M(J/\psi)-M(\eta_c)]/4$, $v^m_{b\bar b}=[M(\Upsilon)-M(\eta_b)]/4$,
$v^m_{b\bar c}$ interpolated between them with $|\psi(0)|^2\propto\mu^p$ in the reduced mass $\mu$~\cite[Sec.~VII]{karliner2014}, and
the same $C$ for a heavy--heavy quark pair. The binding terms are
\begin{equation}\label{eq:exo-B}
B_{ij}=B^{(1)}_{ij}\,\frac{\langle\boldsymbol\lambda_i\!\cdot\!\boldsymbol\lambda_j\rangle}{-16/3},
\end{equation}
with colour-singlet values $B^{(1)}$ for $(c,\bar s)$, $(b,\bar s)$, $(c,\bar c)$, $(b,\bar b)$ and $(b,\bar c)$ fitted to mesons, and none
for pairs with a $u$ or $d$ quark. In a baryon or diquark this is KR's rule $B(QQ)=B(Q\bar Q)/2$; for a pair in different halves of a
multiquark it gives $B^{(1)}/4$, KR's stated assumption that the potential factorises into a colour factor and one spatial factor,
applied to every pair.

\paragraph{Variants.} Two variants were fixed with the model and are reported for every state, never scored. \emph{P-in} drops the
binding between constituents of different halves. \emph{J} replaces the two mass sets by one universal set and adds a constant $S$
per string junction~\cite{karliner2017,karliner2017a}: none for a meson, one for a baryon, two for a tetraquark, three for a pentaquark.
In the tree these are the internal nodes below the root.

\paragraph{Fit, uncertainty and scoring.} All parameters are fitted by unweighted least squares to 42 ordinary ground-state hadrons:
$\pi,\rho,\omega,K,K^*,\phi$, $D^{(*)},D_s^{(*)},B^{(*)},B_s^{(*)},\eta_c,J/\psi,\eta_b,\Upsilon(1S),B_c$, the light octet and decuplet
baryons, and the singly charmed and bottom baryons $\Lambda_Q,\Sigma_Q^{(*)},\Xi_Q,\Xi_Q',\Xi_Q^*,\Omega_Q^{(*)}$ (no $\Omega_b^*$). The
masses are the values tabulated by KR~\cite{karliner2014} or current values of the Particle Data Group~\cite{navas2024}, read state by
state; the $\Xi_b'$ uses LHCb's mass difference~\cite{lhcb2015h}. Isospin averages are taken, since the model has no isospin breaking;
$\eta$ and $\eta'$ are left out, as by KR, because of flavour mixing. Nothing exotic enters the fit, and the doubly charmed
$\Xi_{cc}^{++}$ is held out: it tests the rule $B(cc)=B(c\bar c)/2$ that every doubly heavy prediction uses. The model uncertainty is
$\sigma_{\rm model}=\sqrt{\mathrm{RSS}/(N-p)}$, and a prediction carries $\sigma_{\rm pred}^2=\sigma_{\rm model}^2+g^{\mathsf T}\Sigma g$,
with $g$ the gradient of the predicted mass, $\Sigma=\sigma_{\rm model}^2(\mathsf D^{\mathsf T}\mathsf D)^{-1}$ the parameter covariance
and $\mathsf D$ the Jacobian of the 42 fitted masses with respect to the $p=18$ parameters. A confrontation is scored by
\begin{equation}\label{eq:exo-pull}
z=\frac{M_{\rm model}-M_{\rm data}}{\sqrt{\sigma_{\rm pred}^2+\sigma_{\rm data}^2}},
\end{equation}
with the data error on the side of the model; it \emph{holds} for $|z|<2$ and is \emph{refuted} otherwise, as in Section~\ref{sec:data}.
States not yet seen are \emph{live}, with their mass relative to the lowest two-hadron threshold allowed in an S wave. Every assignment
of a model level to a measured state was fixed before the comparison (Table~\ref{tab:exo-a}).

\paragraph{Hosting check.} With KR's published parameters the implementation reproduces KR's tables to their quoted precision
(0.15\,MeV for the heavy hadrons, the rounding to 1\,MeV for the light ones): the light hadrons, the charmed and bottom baryons, the
$\Xi_{cc}$ at 3626.7\,MeV, and the tetraquarks $T(bb\bar u\bar d)=10389.3$, $T(cc\bar u\bar d)=3882.15$ and
$T(bc\bar u\bar d)=7133.75$\,MeV against their 10389.4, 3882.2 and 7133.7\,MeV~\cite{karliner2017}. It also reproduces
the fully charmed $0^{++}$ estimate of Karliner, Nussinov and Rosner, 6191.5\,MeV~\cite{karliner2017a}, from their formula
\chk{exotics:H}.

\paragraph{The fit.} The fit gives $\sigma_{\rm model}=8.1$\,MeV for 24 degrees of freedom, root-mean-square residuals of 5.5\,MeV for
mesons and 6.6\,MeV for baryons, and a largest residual of $-16.3$\,MeV for the $\Sigma$ (Fig.~\ref{fig:exo-fit}). Fitting with each
hadron left out in turn and predicting it gives 11.2\,MeV for 37 hadrons (the five that fix a parameter exactly are skipped). Our
parameters (Table~\ref{tab:exo-par}) differ from KR's mainly in $m^b_c$ (1715.9 against 1710.5\,MeV) and $B^{(1)}_{c\bar c}$ ($-247.5$
against $-258.0$\,MeV), because KR fix each parameter on one or two hadrons and we fit all at once with current charmonium masses. The
held-out $\Xi_{cc}^{++}$ comes out at $3642.9\pm12.3$\,MeV against LHCb's $3621.55\pm0.38$\,MeV~\cite{lhcb2020j}: $z=+1.74$, the rule
holds \chk{exotics:F}. Variant J fits as well, $\sigma_{\rm model}=8.2$\,MeV with $S=160.6$\,MeV; the ordinary hadrons cannot tell the
two mass conventions apart.

\begin{table}[tbp]\centering\small
\caption{Fitted parameters (1$\sigma$ from the covariance scaled by $\sigma_{\rm model}$) and KR's values~\cite{karliner2014}, all in
MeV; the rows $b$ and $a$ give $b/(310\,\text{MeV})^2$ and $a/(363\,\text{MeV})^2$ \chk{exotics:F}.}
\label{tab:exo-par}
\begin{tabular}{@{}lrrlrr@{}}
\toprule
parameter & fit & KR & parameter & fit & KR \\
\midrule
$m^m_q$ & $309.6\pm2.1$ & 310 & $m^b_q$ & $364.0\pm1.4$ & 363 \\
$m^m_s$ & $481.8\pm3.1$ & 483 & $m^b_s$ & $536.0\pm1.8$ & 538 \\
$m^m_c$ & $1658.1\pm5.2$ & 1663.3 & $m^b_c$ & $1715.9\pm4.4$ & 1710.5 \\
$m^m_b$ & $5001.6\pm5.4$ & 5003.8 & $m^b_b$ & $5047.8\pm4.5$ & 5043.5 \\
$b$ & $79.3\pm1.5$ & 80 & $a$ & $51.1\pm1.0$ & 50 \\
$k_s$ & $1.79\pm0.12$ & 1.65 (baryons) & $B^{(1)}_{c\bar s}$ & $-66.1\pm6.7$ & $-69.9$ \\
$B^{(1)}_{b\bar s}$ & $-80.4\pm7.0$ & $-83.6$ & $B^{(1)}_{c\bar c}$ & $-247.5\pm12.3$ & $-258.0$ \\
$B^{(1)}_{b\bar b}$ & $-558.4\pm12.6$ & $-562.8$ & $B^{(1)}_{b\bar c}$ & $-334.6\pm12.4$ & $-341.5$ \\
$v^m_{c\bar c}$ & $28.2\pm2.9$ & 28.3 & $v^m_{b\bar b}$ & $15.4\pm2.9$ & 15.6 \\
\bottomrule
\end{tabular}
\end{table}

\begin{figure}[tbp]\centering
\includegraphics[width=0.49\linewidth]{exotic_fit.pdf}
\caption{Fit of the colour-spin model to 42 ordinary ground-state hadrons: model minus measured mass against mass, mesons (circles) and
baryons (squares); open circles are the residuals when the hadron is left out of the fit; the band is $\pm\sigma_{\rm model}=\pm8.1$\,MeV.
The diamond is the held-out $\Xi_{cc}^{++}$ with $\pm\sigma_{\rm pred}$. Data: KR's tables~\cite{karliner2014} and the Particle Data
Group~\cite{navas2024}. What to look for: no trend with mass, and the held-out baryon within $2\sigma$ \chk{exotics:F}.}
\label{fig:exo-fit}
\end{figure}

\subsection{Compact exotics against data}\label{sec:exotics-compact}

Table~\ref{tab:exo-a} lists the model's compact trees and the measurements they are confronted with, Fig.~\ref{fig:exo-spectrum} shows
the pulls and the doubly heavy binding energies, and Fig.~\ref{fig:exo-levels} the pentaquark and fully charmed levels \chk{exotics:A}.

\begin{table}[tbp]\centering\footnotesize\setlength{\tabcolsep}{4pt}
\caption{Compact H8 trees: the model (KR's convention, our fit) with $\sigma_{\rm pred}$, the same tree with KR's own parameters, the
variants P-in and J, the measured mass with its error, the pull~\eqref{eq:exo-pull} and the verdict. The last column is the model mass
minus the lowest S-wave two-hadron threshold. Data: $T_{cc}^+$~\cite{lhcb2022i}, $X_0(2900)$~\cite{lhcb2020i}, $T_{c\bar s0}(2900)$
\cite{lhcb2023h}, $X(6900)$~\cite{lhcb2020h}, $P_c$~\cite{lhcb2019h}, $P_{\psi s}(4338)$~\cite{lhcb2023j}, $P_{cs}(4459)$
\cite{lhcb2021h}, $X(3872)$~\cite{lhcb2020k}, $Z_c(3900)$~\cite{besiii2013h}. $^\dagger$: the level nearest to the datum, an
assignment by proximity; the chance that a mass placed at random within $\pm150$\,MeV lies as close to some level is 0.57,
0.24 and 0.08 for $P_c(4440)$, $P_c(4457)$ and $P_{cs}(4459)$. Masses in MeV \chk{exotics:A}.}
\label{tab:exo-a}
\begin{tabular}{@{}llrrrrrrlr@{}}
\toprule
state & tree, $J^{P(C)}$ & model & KR par. & P-in & J & data & $z$ & verdict & $M-T$ \\
\midrule
$T_{cc}^+$ & $[cc]_1[\bar u\bar d]_0$, $1^+$ & $3897\pm12$ & 3882 & 3897 & 3995 & $3874.74\pm0.08$ & $+1.86$ & holds & $+22$ \\
$X_0(2900)$ & $[ud]_0[\bar c\bar s]_0$, $0^+$ & $2755\pm10$ & 2754 & 2755 & 2861 & $2866\pm7$ & $-9.28$ & refuted & $+392$ \\
$T_{c\bar s0}(2900)$ & $[cq][\bar s\bar q]$, $0^+$ & $2809\pm9$ & 2812 & 2825 & 2916 & $2908\pm23$ & $-4.06$ & refuted & $+701$ \\
$X(6900)$ & $[cc]_1[\bar c\bar c]_1$, $2^{++}$ & $6425\pm29$ & 6383 & 6672 & 6515 & $6905\pm13$ & $-15.2$ & refuted & $+231$ \\
$P_c(4312)$ & lowest $([cq][qq])\bar c$ & $4277\pm11$ & 4264 & 4339 & 4485 & $4311.9^{+6.8}_{-0.9}$ & $-3.06$ & refuted & $+354$ \\
$P_c(4440)^\dagger$ & $([cq][qq]_1)\bar c$, $\frac12^-$ & $4428\pm11$ & 4446 & 4396 & 4485 & $4440.3^{+4.3}_{-4.9}$ & $-0.98$ & holds & $+506$ \\
$P_c(4457)^\dagger$ & $([cq][qq]_1)\bar c$, $\frac12^-$ & $4462\pm11$ & 4456 & 4490 & 4485 & $4457.3^{+4.1}_{-1.8}$ & $+0.38$ & holds & $+539$ \\
$P_{\psi s}(4338)$ & $([cs][ud]_0)\bar c$, $\frac12^-$ & $4392\pm11$ & 4382 & 4470 & 4602 & $4338.2\pm0.8$ & $+4.81$ & refuted & $+292$ \\
$P_{cs}(4459)^\dagger$ & $([cs][ud]_0)\bar c$, $\frac32^-$ & $4460\pm11$ & 4442 & 4470 & 4602 & $4458.8^{+5.5}_{-3.1}$ & $+0.13$ & holds & $+361$ \\
\midrule
\multicolumn{10}{@{}l}{\emph{not scored}}\\
$X(6900)$ & $[cc]_1[\bar c\bar c]_1$, $0^{++}$ & $6340\pm26$ & 6298 & 6588 & 6430 & $6905\pm13$ & $-19.1$ & info & $+372$ \\
$X(3872)$ & $[cq][\bar c\bar q]$, $1^{++}$ & $4108\pm11$ & 4093 & 4170 & 4212 & $3871.70\pm0.10$ & $+21.0$ & info & $+237$ \\
$Z_c(3900)$ & $[cq][\bar c\bar q]$, $1^{+-}$ & $3985\pm11$ & 3969 & 4047 & 4087 & $3899\pm6$ & $+6.72$ & info & $+749$ \\
\midrule
\multicolumn{10}{@{}l}{\emph{live}}\\
$T_{bb}^-$ & $[bb]_1[\bar u\bar d]_0$, $1^+$ & $10399\pm12$ & 10389 & 10399 & 10512 & -- & -- & live & $-205$ \\
$T_{bc}$ & $[bc]_0[\bar u\bar d]_0$, $0^+$ & $7146\pm11$ & 7134 & 7146 & 7251 & -- & -- & live & $+2$ \\
$T_{bc}$ & $[bc]_1[\bar u\bar d]_0$, $1^+$ & $7180\pm12$ & 7168 & 7180 & 7285 & -- & -- & live & $-9$ \\
$P_{cc}$ & $([cc]_1[ud]_0)\bar q$, $\frac12^-$ & $4232\pm12$ & 4215 & 4232 & 4436 & -- & -- & live & $+471$ \\
$cc\bar c\bar c$ & $[cc]_1[\bar c\bar c]_1$, $1^{+-}$ & $6369\pm27$ & 6326 & 6616 & 6459 & -- & -- & live & $+288$ \\
\bottomrule
\end{tabular}
\end{table}

\begin{figure}[tbp]\centering
\includegraphics[width=\linewidth]{exotic_spectrum.pdf}
\caption{Left: model minus measured mass for the nine scored confrontations of Table~\ref{tab:exo-a} (symmetric logarithmic axis,
linear within $\pm30$\,MeV), for the model with its $\pm1\sigma$ combined error and the variants P-in and J; $\dagger$ marks the
assignments by proximity. Right: mass minus the lowest S-wave threshold for the doubly heavy tetraquarks: the model, variant J, KR's
published values~\cite{karliner2017}, LHCb's $T_{cc}^+$~\cite{lhcb2022i}, and lattice QCD~\cite{francis2017,leskovec2019,aoki2023,
hudspith2023,tripathy2025,francis2019,padmanath2024}. What to look for: the scored points scatter far beyond $\pm2\sigma$ except
$T_{cc}^+$ and the proximity assignments; no convention places them all \chk{exotics:A}.}
\label{fig:exo-spectrum}
\end{figure}

\paragraph{Doubly heavy tetraquarks.} $T_{cc}^+$ is LHCb's narrow $D^0D^0\pi^+$ state, with pole mass $360\pm40$\,keV below the
$D^{*+}D^0$ threshold (3875.10\,MeV), $J^P=1^+$ and isospin zero~\cite{lhcb2022h,lhcb2022i}. The tree $[cc]_1[\bar u\bar d]_0$ gives
$3897\pm12$\,MeV, $z=+1.86$: \emph{holds}. The model places it 22\,MeV above the threshold where it is bound by 0.36\,MeV; with
$\sigma_{\rm pred}=12$\,MeV it cannot decide whether the state is bound, as KR said of their own 3882\,MeV~\cite{karliner2017}. The
same tree predicts $T_{bb}^-$ at $10399\pm12$\,MeV, 205\,MeV below the $B^-\bar B^{*0}$ threshold and stable against strong and
electromagnetic decay (KR: 215\,MeV). Lattice QCD finds it bound too, mostly by less and with a large spread: $189\pm10$~\cite{francis2017},
$128\pm26$~\cite{leskovec2019}, $83\pm22$~\cite{aoki2023}, $112\pm13$~\cite{hudspith2023} and $116^{+36}_{-30}$\,MeV~\cite{tripathy2025}
(binding energies relative to $B\bar B^*$), pulls from $-1.0$ to $-5.1$ against our value (these are calculations, not
measurements, and enter no verdict). For $bc\bar u\bar d$ the model gives a $0^+$ state 2\,MeV above $\bar B^0D^0$ and a $1^+$ state 9\,MeV below $D^0\bar B^{*}$, both within one $\sigma_{\rm pred}$
of threshold; lattice QCD finds the $1^+$ bound by 15 to 61\,MeV~\cite{francis2019} or by $\Delta E=-43^{+15}_{-25}$\,MeV
\cite{padmanath2024}, pulls $+1.1$ and $+1.8$, and Alexandrou et al.\ find shallow bound or virtual states a few MeV or less below both
thresholds~\cite{alexandrou2024}.
\emph{Verdict: $T_{bb}$ and $T_{bc}$ live.} Their isospin-one partners are in Section~\ref{sec:exotics-partners}.

\paragraph{Open charm and strangeness.} LHCb's $X_0(2900)$, a spin-0 $\bar cd\bar su$ state at $2866\pm7$\,MeV in $B^+\to D^+D^-K^+$
\cite{lhcb2020i}, would be the tree $[ud]_0[\bar c\bar s]_0$; the model puts it at $2755\pm10$\,MeV, $z=-9.3$: \emph{refuted}. The
$T_{c\bar s0}(2900)^{++,0}$ seen in $D_s^+\pi^\pm$, an isospin triplet with $J^P=0^+$ at $2908\pm23$\,MeV~\cite{lhcb2023h}, would be the
lowest $0^+$ of $[cq][\bar s\bar q]$; the model gives $2809\pm9$\,MeV, $z=-4.1$: \emph{refuted}. The spin-1 $X_1(2900)$ has negative
parity, so it needs an orbital excitation and lies outside the S-wave model. The two refutations are where the mass conventions matter
most: variant J would put both states within $0.5\sigma$ of the data, at the price of $T_{cc}^+$, which it places 120\,MeV too high
($z=+10.8$).

\paragraph{Four charm quarks.} The S-wave tree $[cc]_1[\bar c\bar c]_1$ gives $0^{++}$, $1^{+-}$ and $2^{++}$ at 6340, 6369 and
6425\,MeV, 146 to 231\,MeV above the $J/\psi J/\psi$ threshold. LHCb's narrow $X(6900)$, $6905\pm11\pm7$\,MeV (no interference with the
continuum)~\cite{lhcb2020h}, lies 480\,MeV above even the most favourable assignment, the $2^{++}$: $z=-15$, \emph{refuted}; it is
refuted in both variants as well. The structures CMS and ATLAS report~\cite{cms2024h,atlas2023h} are shown in Fig.~\ref{fig:exo-levels}
but not scored: CMS's no-interference fit finds three at $6552\pm16$, $6927\pm10$ and $7287^{+21}_{-19}$\,MeV, ATLAS's model~A three at
$6.41^{+0.11}_{-0.09}$, $6.63^{+0.09}_{-0.05}$ and $6.86^{+0.03}_{-0.04}$\,GeV, and LHCb a broad enhancement from 6.2 to 6.8\,GeV.
Only the lowest ATLAS structure lies near the S-wave levels; that was noticed after the comparison and is not a test. The narrow
states near 6.9\,GeV would have to be excitations, which this model does not compute.

\begin{figure}[tbp]\centering
\includegraphics[width=\linewidth]{exotic_levels.pdf}
\caption{Levels of the compact trees (model, KR's convention, our fit). Left: hidden-charm pentaquarks $([cq][qq])\bar c$ with $I=1/2$
by $J$ (all of negative parity), coloured by the spin of the light diquark, against LHCb's $P_c(4312)$, $P_c(4440)$,
$P_c(4457)$~\cite{lhcb2019h} (solid) and the $\Sigma_c\bar D^{(*)}$, $\Sigma_c^*\bar D$ thresholds (dashed). Centre: $udsc\bar c$ with
$I=0$ in the two trees $([cs][ud]_0)\bar c$ and $([cq][qs])\bar c$ against $P_{\psi s}(4338)$~\cite{lhcb2023j} and $P_{cs}(4459)$
\cite{lhcb2021h}, with the $\Xi_c\bar D^{(*)}$ thresholds. Right: $[cc]_1[\bar c\bar c]_1$ for the model and both variants against the
di-$J/\psi$ structures of LHCb (model I)~\cite{lhcb2020h}, CMS (no interference)~\cite{cms2024h} and ATLAS (model A)~\cite{atlas2023h}.
What to look for: the lowest model levels sit 35 to 55\,MeV away from the lowest observed pentaquarks, the levels are dense enough above
4.4\,GeV to match any peak, and the S-wave four-charm levels lie far below 6.9\,GeV \chk{exotics:A}.}
\label{fig:exo-levels}
\end{figure}

\paragraph{Pentaquarks.} Maiani, Polosa and Riquer read the hidden-charm pentaquarks as $\bar c[cq][q'q'']$, three colour-$\mathbf{\bar 3}$
units~\cite{maiani2015}; in H8's form, $([cq][qq])\bar c$. The S-wave multiplet has 10 levels with isospin $1/2$, from 4277 to
4584\,MeV, with the multiplicities of Ref.~\cite{maiani2015} \chk{exotics:E}. LHCb sees three narrow peaks in $J/\psi p$ with
undetermined $J^P$~\cite{lhcb2019h}. The pre-registered test compares the lowest level, $4277\pm11$\,MeV (a spin-0 light diquark,
$J=1/2$), with the lowest peak, $P_c(4312)$: $z=-3.1$, \emph{refuted} (variant P-in, without the $c\bar c$ binding across the halves,
gives $z=+2.0$). $P_c(4440)$ and $P_c(4457)$ match the nearest levels ($z=-1.0$ and $+0.4$), but with 10 levels a mass placed at random
within $\pm150$\,MeV would match as well with probability 0.57 and 0.24: these two hold and test little. For the strange pentaquark
$P_{\psi s}(4338)$, with $J=1/2$ and preferred negative parity~\cite{lhcb2023j}, the lowest $J=1/2$ level with $I=0$ is
$([cs][ud]_0)\bar c$ at $4392\pm11$\,MeV: $z=+4.8$, \emph{refuted}. The evidence-level $P_{cs}(4459)$~\cite{lhcb2021h} lies 1.5\,MeV from
a $J=3/2$ level, a chance coincidence with probability 0.08. The doubly charmed pentaquark $([cc]_1[ud]_0)\bar q$ is a live prediction at
$4232\pm12$\,MeV, 471\,MeV above $\Xi_{cc}\pi$, into which it would decay strongly.

\paragraph{Hidden charm.} For the record, the compact $[cq][\bar c\bar q]$ trees give $1^{++}$ at $4108\pm11$\,MeV and $1^{+-}$ at
$3985\pm11$\,MeV, 236 and 86\,MeV above $X(3872)$ and $Z_c(3900)$~\cite{lhcb2020k,besiii2013h}. In this model $X(3872)$ is not a compact
tree; it sits on the $D^0\bar D^{*0}$ threshold, as a molecular tree would.

\paragraph{Summary of (a).} Of nine scored confrontations, four hold and five are refuted; three of the four that hold are assignments
by proximity, and the fourth, $T_{cc}^+$, holds with $z=+1.86$. Each mass convention rescues a different subset: the model holds 4,
P-in 3 and J 2 of the nine, and no convention fits $T_{cc}^+$, $X_0(2900)$ and the pentaquarks together. \emph{Verdict: refuted} as a
model of the compact exotics; $T_{bb}$, $T_{bc}$ and $P_{cc}$ \emph{live}.

\subsection{Threshold proximity}\label{sec:exotics-thresholds}

Where both halves are colour singlets, H8 makes the state a molecule of two hadrons, and a loosely bound molecule lies just below the
threshold of its constituents. We test whether the established exotic hadrons lie closer to two-hadron thresholds than chance would
place them.

\paragraph{Catalogue and thresholds.} The catalogue (Table~\ref{tab:exo-b}) holds the 26 exotic hadrons reported as observations,
not evidence, with a measured mass, each from its primary publication: $\chi_{c1}(3872)$~\cite{lhcb2020k,lhcb2013h},
$Z_c(3900)$~\cite{besiii2013h,besiii2017h}, $Z_c(4020)$~\cite{besiii2013i}, $Z_c(4430)$~\cite{lhcb2014h}, $Z_{cs}(3985)$
\cite{besiii2021h}, $Z_{cs}(4000)$, $Z_{cs}(4220)$, $\chi_{c1}(4140)$, $\chi_{c1}(4274)$, $\chi_{c0}(4500)$, $\chi_{c0}(4700)$, $X(4685)$,
$X(4630)$~\cite{lhcb2021i,lhcb2017h}, $X(3960)$~\cite{lhcb2023i}, $\psi(4230)$~\cite{besiii2017i}, $Z_b(10610)$, $Z_b(10650)$
\cite{belle2012h}, $T_{cc}^+$~\cite{lhcb2022i}, $X_0(2900)$, $X_1(2900)$~\cite{lhcb2020i}, $T_{c\bar s0}(2900)$~\cite{lhcb2023h},
$P_c(4312)$, $P_c(4440)$, $P_c(4457)$~\cite{lhcb2019h}, $P_{\psi s}(4338)$~\cite{lhcb2023j} and $X(6900)$~\cite{lhcb2020h}. For each we
list every pair of hadrons narrower than 60\,MeV (83 hadrons and their antiparticles: light, open-charm and open-bottom mesons including
the narrow P-wave ones, charmonia and bottomonia below open flavour, light, charmed and bottom baryons~\cite{navas2024}) whose
quark flavours add up to the state's, and keep the pairs that can form it in an S wave: $J$ in $|j_1-j_2|,\dots,j_1+j_2$, parity
$P_1P_2$ and, for neutral states with hidden flavour, $C=C_1C_2$ for two self-conjugate hadrons or $(-1)^S$ for a particle and its own
antiparticle. Unknown quantum numbers impose nothing; isospin and $G$ parity are not imposed. The nearest allowed threshold $T$ gives
the distance $d=|M-T|$.

\paragraph{Null and statistic.} Under the null each mass is placed uniformly in $[M-W,M+W]$ with $W=150$\,MeV, and
\begin{equation}\label{eq:exo-p}
p_i=\frac{1}{2W}\Big|\,[M_i-W,M_i+W]\cap\bigcup_k[T_k-d_i,T_k+d_i]\Big|
\end{equation}
is the probability that the nearest allowed threshold is at most as far as observed, computed exactly from the threshold list (it agrees
with 20000 random placements per state). The $p_i$ are combined by Fisher's method, $X^2=-2\sum_i\ln p_i\sim\chi^2_{52}$. The windows
contain 5.5 allowed thresholds on average; this is the look-elsewhere count, and it is built into each $p_i$.

\paragraph{Result.} Thirteen of the 26 states lie within 10\,MeV of an allowed threshold, where 7.0 are expected (Poisson $p=0.027$),
and 7 within 5\,MeV where 3.8 are expected. Fisher's statistic is $X^2=101.1$, $p=5.4\times10^{-5}$; with $W=300$\,MeV, $p=4.1\times
10^{-6}$ (Fig.~\ref{fig:exo-thresholds}) \chk{exotics:B}. The result rests on few states: $X(3872)$ sits 0.005\,MeV above
$D^0\bar D^{*0}$ ($p_i=10^{-4}$) and $T_{cc}^+$ 0.36\,MeV below $D^{*+}D^0$ ($p_i=0.007$); without these two, $p=0.057$. When every
distance is floored at the state's mass error, since a smaller distance is not resolved, $p=5.1\times10^{-3}$. Restricted to the 17 states
with measured $J^P$, $p=1.8\times10^{-4}$. Perturbing the 45 threshold masses we could not re-read by 1 (3)\,MeV moves $p$ within
$7.0\times10^{-6}$ to $7.5\times10^{-5}$ ($5.0\times10^{-7}$ to $1.4\times10^{-4}$). Twelve of the 26 lie below their nearest threshold
and fourteen above.

\begin{table}[tbp]\centering\scriptsize\setlength{\tabcolsep}{2.6pt}
\caption{Threshold proximity: mass (MeV), $J^{P(C)}$ (? unknown), nearest allowed S-wave threshold, $M-T$ (MeV) and $p_i$ of
Eq.~\eqref{eq:exo-p} for $W=150$\,MeV. Sources in the text \chk{exotics:B}.}
\label{tab:exo-b}
\begin{tabular}{@{}lrlrrr@{\hspace{8pt}}lrlrrr@{}}
\toprule
state & $M$ & $J^{P(C)}$ & threshold & $M-T$ & $p_i$ & state & $M$ & $J^{P(C)}$ & threshold & $M-T$ & $p_i$ \\
\midrule
$\chi_{c1}(3872)$ & 3871.7 & $1^{++}$ & $D^0\bar D^{*0}$ & $+0.0$ & 0.0001 & $\psi(4230)$ & 4222.0 & $1^{--}$ & $\omega\chi_{c0}$ & $+24.6$ & 0.59 \\
$Z_c(3900)$ & 3899.0 & $1^{+}$ & $\psi(3770)\pi$ & $-14.3$ & 0.39 & $Z_b(10610)$ & 10607.2 & ? & $B\bar B^*$ & $+2.8$ & 0.075 \\
$Z_c(4020)$ & 4022.9 & ? & $D^*\bar D^*$ & $+5.8$ & 0.11 & $Z_b(10650)$ & 10652.2 & ? & $B^*\bar B^*$ & $+2.7$ & 0.055 \\
$Z_c(4430)$ & 4475.0 & $1^{+}$ & $D_1\bar D_1$ & $-373.3$ & 0.72 & $T_{cc}^+$ & 3874.7 & $1^+$ & $D^0D^{*+}$ & $-0.4$ & 0.0072 \\
$Z_{cs}(3985)$ & 3982.5 & ? & $\bar D^0D_s^{*+}$ & $+5.5$ & 0.34 & $X_0(2900)$ & 2866.0 & $0^+$ & $D^{*-}K^{*+}$ & $-36.1$ & 0.24 \\
$Z_{cs}(4000)$ & 4003.0 & $1^{+}$ & $J/\psi K^{*+}$ & $+14.2$ & 0.33 & $X_1(2900)$ & 2904.0 & $1^-$ & $\bar D_1^0K^0$ & $-15.7$ & 0.11 \\
$Z_{cs}(4220)$ & 4216.0 & $1^{?}$ & $\psi(2S)K^+$ & $+36.2$ & 0.89 & $T_{c\bar s0}(2900)$ & 2908.0 & $0^+$ & $D^{*+}K^{*+}$ & $+5.9$ & 0.039 \\
$\chi_{c1}(4140)$ & 4118.0 & $1^{++}$ & $J/\psi\phi$ & $+1.6$ & 0.022 & $P_c(4312)$ & 4311.9 & ? & $\Sigma_c^+\bar D^0$ & $-5.6$ & 0.25 \\
$\chi_{c1}(4274)$ & 4294.0 & $1^{++}$ & $\omega\psi(2S)$ & $-174.8$ & 0.99 & $P_c(4440)$ & 4440.3 & ? & $p\chi_{c1}$ & $-8.6$ & 0.55 \\
$\chi_{c0}(4500)$ & 4474.0 & $0^{++}$ & $\omega\psi(2S)$ & $+5.2$ & 0.10 & $P_c(4457)$ & 4457.3 & ? & $\Lambda_c(2595)\bar D^0$ & $+0.2$ & 0.022 \\
$\chi_{c0}(4700)$ & 4694.0 & $0^{++}$ & $\phi\psi(2S)$ & $-11.6$ & 0.42 & $P_{\psi s}(4338)$ & 4338.2 & $\frac12^-$ & $\Xi_c^+D^-$ & $+0.8$ & 0.063 \\
$X(4685)$ & 4684.0 & $1^{++}$ & $\phi\psi(2S)$ & $-21.6$ & 0.49 & $X(6900)$ & 6905.0 & $?^{?+}$ & $\chi_{c0}\chi_{c1}$ & $-20.4$ & 0.86 \\
$X(4630)$ & 4626.0 & $1^{-+}$ & $D_s^*\bar D_{s1}(2536)$ & $-21.3$ & 0.79 & $X(3960)$ & 3956.0 & $0^{++}$ & $\eta'\eta_c$ & $+14.1$ & 0.33 \\
\bottomrule
\end{tabular}
\end{table}

\begin{figure}[tbp]\centering
\includegraphics[width=\linewidth]{exotic_thresholds.pdf}
\caption{Threshold proximity of the 26 exotic hadrons of Table~\ref{tab:exo-b}. Left: distance to the nearest allowed S-wave
two-hadron threshold (bars, logarithmic) and the median distance under the null of uniform placement within $\pm150$\,MeV (ticks).
Right: empirical distribution of the per-state $p_i$ against the uniform null, for $W=150$ and 300\,MeV, with Fisher's combined $p$. What
to look for: the excess of small $p_i$ comes from a handful of states, two of them within 1\,MeV of threshold \chk{exotics:B}.}
\label{fig:exo-thresholds}
\end{figure}

\paragraph{Reading.} Exotic hadrons do lie preferentially near thresholds. That is what the molecular branch of H8 says, but it is
also the standard phenomenology of hadronic molecules and of threshold cusps, which QCD produces without any reference to GUID, and it
does not choose between them. The catalogue is not a random sample: many of these states were found as narrow structures close to a
threshold, where a cusp can mimic a peak, so part of the excess is selection. And the significance is carried by two states.
\emph{Verdict: holds}, as a statement about exotic hadrons; as evidence for H8 over QCD, \emph{holds as a reading}.

\subsection{Partners}\label{sec:exotics-partners}

A tree with light-quark halves comes in isospin multiplets, and a tree whose halves can carry spin comes in spin multiplets. We compare
each predicted partner with the searches we could verify \chk{exotics:C}.

\begin{enumerate}[leftmargin=*,itemsep=2pt]
\item \emph{Charged partners of a compact $X(3872)$.} The diquark--antidiquark reading of $X(3872)$ of Maiani et al.~\cite{maiani2005}
has charged partners $X^\pm$ near 3872\,MeV. BaBar searched for $X^-\to J/\psi\pi^-\pi^0$ in $B\to X^-K$ and set
$\mathcal B(B^0\to X^-K^+,X^-\to J/\psi\pi^-\pi^0)<5.4\times10^{-6}$ at 90\% CL where an isovector partner implies $2.56\times10^{-5}$,
75 signal events against none seen~\cite{babar2005h}; Belle's limits lie well below the isospin-triplet expectation as
well~\cite{belle2011h}. \emph{Failed prediction} of the compact reading. Our model agrees that $X(3872)$ is not compact: it puts the
compact $1^{++}$ isomultiplet, charged members included, at $4108\pm11$\,MeV, where no search was verified (\emph{untested}).
\item \emph{Isospin partners of $T_{cc}^+$.} The tree $[cc]_1[\bar q\bar q]_1$ with $I=1$ gives $J=0,1,2$ at 4042, 4071 and
4131\,MeV, 312, 200 and 117\,MeV above their lowest S-wave thresholds, so broad. LHCb sees no narrow structure in $D^+D^0\pi^+$ (3.9 to
4.0\,GeV) or $D^+D^+$ (up to 3.87\,GeV) and disfavours an isovector $T_{cc}^+$~\cite{lhcb2022i}, as the model has it; the predicted
partners lie outside those windows: \emph{untested}. A good antidiquark has spin 0, so the tree has no isoscalar spin partner.
\item \emph{Isospin partners of $T_{bb}^-$}: 10583, 10593 and 10613\,MeV for $J=0,1,2$, $+25$, $-11$ and $-37$\,MeV from their lowest
thresholds: \emph{live}.
\item \emph{Pentaquark multiplets.} The 10 levels with $I=1/2$ have partners with $I=3/2$, degenerate with the 7 levels that contain a
spin-1 light diquark. LHCb's $\Lambda_b\to J/\psi pK^-$ shows three narrow peaks, and production in $\Lambda_b$ decays favours a spin-0
$[ud]$~\cite{maiani2015}, so missing peaks exclude nothing; no search for $I=3/2$ hidden-charm pentaquarks was verified:
\emph{untested}.
\item \emph{Four-charm and open-strange partners.} The $1^{+-}$ of $[cc]_1[\bar c\bar c]_1$ cannot decay to $J/\psi J/\psi$, and the
$1^+$ partner of $X_0(2900)$, $[ud]_0[\bar c\bar s]_1$ at 2807\,MeV, cannot decay to $D^-K^+$: \emph{untested}.
\end{enumerate}

\noindent One predicted partner is excluded by a search: the charged $X^\pm$ of the compact $X(3872)$. \emph{Verdict: refuted} for the
compact reading of $X(3872)$; every other partner is untested or live.

\paragraph{Choices made after seeing the data.} (i) The binding rule across the halves (primary: every pair; variant P-in: none) was
written into the script before any exotic state was computed, but by authors who knew the published masses; both are reported for
every state. (ii) Sharing $k_s$ between mesons and baryons was fixed before the fit, from KR's published $D_s$ residuals, not from any
exotic. (iii) For each catalogue state we took one measurement, the most recent one we could read from the cited collaboration (LHCb's
2021 values for the $J/\psi\phi$ states; BESIII's discovery values for $Z_c(3900)$ and $Z_c(4020)$, not the averages of the Particle Data
Group, which we could not read); we left out $X(3915)$ (its exotic nature is disputed), $P_{cs}(4459)$ and $P_c(4337)$ (evidence below
5 standard deviations). (iv) The chance probabilities of the proximity assignments were computed after seeing that they hold. (v) The
tests without the two sub-MeV states, with the measured-$J^P$ subset and with distances floored at the mass error were added after the
Fisher result, to show what carries it. (vi) The remark on the lowest ATLAS structure was made after the comparison.

\paragraph{Caveats.} The mass convention for multiquarks dominates: KR's baryon masses and the junction variant differ by 90 to
113\,MeV for the tetraquarks and 204 to 210\,MeV for the lowest pentaquarks, ten to twenty-five times $\sigma_{\rm model}$, and the
ordinary hadrons cannot choose between them. $\sigma_{\rm model}$ measures ordinary hadrons; carrying the model to four and five
constituents adds uncertainty it does not contain (the size of a diquark, binding across the halves, the colour configurations outside
the tree, such as $\mathbf6\otimes\mathbf{\bar6}$, whose mixing would lower the lowest levels). Only S-wave ground states are computed,
so a mismatch refutes an S-wave assignment, not every compact reading. The threshold test inherits the selection of its catalogue, and
its look-elsewhere count depends on the threshold list. Lattice results for $T_{bb}$ scatter by 100\,MeV. We could not reach the server
of the Particle Data Group; its current values were read through a search of its listings, and 45 of the 83 threshold hadrons were not
re-read (they move no conclusion, see above).

All numbers in this section come from \texttt{docs/simulation/exotic\_hadrons.py} (75 of 75 checks, about one minute on one core).

\paragraph{Verdict.} H8 holds as a reading: QCD's colour algebra already writes every hadron as a binary tree of conjugate halves, and
the compact and molecular trees together span the colour states of a tetraquark, so the hypothesis cannot fail on any single state. Its
content lies in the masses of the compact trees, and there the standard colour-spin model, fitted to ordinary hadrons
($\sigma_{\rm model}=8.1$\,MeV), holds for $T_{cc}^+$ ($z=+1.9$) and is refuted for $X_0(2900)$, $T_{c\bar s0}(2900)$, $X(6900)$,
$P_c(4312)$ and $P_{\psi s}(4338)$; the charged partners of a compact $X(3872)$ are excluded. Exotic hadrons lie closer to two-hadron
thresholds than chance ($p=5\times10^{-5}$, carried by $X(3872)$ and $T_{cc}^+$), as the molecular branch and ordinary QCD molecules
both say. The live predictions are a $T_{bb}^-$ 205\,MeV below $B\bar B^*$ at $10399\pm12$\,MeV and $T_{bc}$ states within 10\,MeV of
their thresholds.
